\documentclass[11pt]{article}

\usepackage[margin=1.0in]{geometry}
\usepackage{amsmath,amssymb}
\usepackage{booktabs}
\usepackage{graphicx}
\usepackage{microtype}
\usepackage{xcolor}
\usepackage[colorlinks=true,linkcolor=blue!60!black,citecolor=blue!60!black,urlcolor=blue!60!black]{hyperref}

\newcommand{\PT}{\mathcal{PT}}
\newcommand{\ee}{\mathrm{e}}
\newcommand{\ii}{\mathrm{i}}

\title{\vspace{-1.2em}Experimental proposal: divisibility onset, generator
floor, and the backflow--divisibility dissociation in a passive $\PT$
photonic coupler\vspace{-0.4em}}
\author{Jonathan Torkelson\\[2pt]
\small Embeddetech, Inc., Tulsa, Oklahoma, USA \quad
\texttt{jtorkelson@embeddetech.com}}
\date{July 11, 2026}

\begin{document}
\maketitle
\vspace{-1.5em}

\noindent\textbf{Summary.} We propose a set of measurements on
femtosecond-laser-written lossy waveguide couplers---the platform of
Klauck \emph{et al.}~\cite{Klauck2019}---that would extract, for the
first time, three process-level observables of the classical
\emph{transition-matrix process} of a $\PT$ dimer: the length $t_c$ at
which divisibility is first lost (permanently, in the broken phase), the asymptotic generator floor
$-(a+\tilde\kappa)$ (gain $a\equiv\gamma/2$, $\tilde\kappa\equiv\sqrt{a^2-\kappa^2}$), and the dissociation between
information-flow Markovianity and divisibility, phase-locked to the
$\PT$ transition. All three follow with
no free parameters from the device constants; classical light
suffices; and the primary observable is exactly invariant under the
per-launch transmission-calibration errors, the leading multiplicative
systematic (though not under the dark floor, which must be
subtracted). A full measurement-chain simulation
(\texttt{scripts/t2\_measurement\_chain.py}, public) underwrites the error
budget quoted below.

\paragraph{Theory targets.} Write the coupler as the dimer
$H=\cos\theta\,\mathbf{1}+\ii\sin\theta\,\sigma_z+s\,\sigma_x$; the
passive device with coupling $\kappa$ [cm$^{-1}$] and single-guide
loss rate $\gamma$ [cm$^{-1}$, amplitude convention, $\PT$-breaking at
$\gamma=2\kappa$~\cite{Klauck2019}] realizes $s\to\kappa$,
$\sin\theta\to a\equiv\gamma/2$, $t\to z$ (propagation length) after
removing the mean loss. The indivisible-process paper~\cite{indivisible}
establishes, for the measured
intensity transfer matrix $\Gamma(z)$ (columns: launch guide; rows:
output guide; index~1 the low-loss guide):
\begin{itemize}\setlength\itemsep{0pt}
\item \emph{Divisibility onset:} P-divisibility first fails at $t_c$,
with
$t_c=\tilde\kappa^{-1}\operatorname{arcsinh}\!\bigl(\tilde\kappa/(\kappa\sqrt2)\bigr)$,
$\tilde\kappa=\sqrt{a^2-\kappa^2}$ (continued to $\arcsin$ in the
unbroken phase), exactly where $\det\Gamma$ crosses
zero; in the broken phase ($\gamma>2\kappa$) the resulting
indivisibility is \emph{permanent} (P-divisible precisely on $[0,t_c)$),
in the unbroken phase it recurs in
windows of period $\pi/\sqrt{\kappa^2-a^2}$, and $t_c$ is continuous
through the exceptional point.
\item \emph{Generator floor:} in the broken phase the off-diagonal
generator entry satisfies $L_{12}(z)<-(a+\tilde\kappa)$ for all
$z>t_c$ and \emph{rises} to the strictly negative constant
$-(a+\tilde\kappa)$ as $z\to\infty$: a floor on $|L_{12}|$,
approached from below by the signed entry.
\item \emph{Dissociation:} on the same data, the survival-conditioned
total-variation distance between the two launch columns shows
\emph{negligible} backflow in the late window (for onsets $t_c$
preceding it) of the broken phase while divisibility fails
terminally---numerically, the two standard
notions of non-Markovianity come
apart, phase-locked to the $\PT$ transition.
\end{itemize}

\paragraph{Why this is new.} The apparatus exists and is
characterized as a function of length~\cite{Klauck2019}; exceptional-point
criticality of information flow has been measured~\cite{Xiao2019}; and
indivisibility without information backflow has already been observed
on one tomographic data set---Bernardes \emph{et al.}'s ``weak''
non-Markovian regime, together with the transition to the strong
one~\cite{Bernardes2015}. What no experiment has extracted, to our
knowledge, is a closed-form divisibility-onset \emph{length} or a
generator \emph{floor}, nor phase-locked the dissociation to a $\PT$
transition on a passive photonic coupler. Note that Klauck
\emph{et al.}\ operated well inside the unbroken phase (far from the
exceptional point at $\gamma/2\kappa=1$); the broken
phase---where the terminal behavior lives---is unexplored territory
for these observables.

\paragraph{Devices and protocol.} Fabricate four couplers (D1--D4) spanning three $\PT$ regimes
(all at $\lambda=815\,$nm in fused silica; loss by sinusoidal
modulation~\cite{Eichelkraut2014}; per sample, one Hermitian reference
coupler, one straight reference guide, and one straight modulated-loss
guide---the latter carrying the device modulation in isolation, so its
length-resolved transmission pins the loss rate $\gamma$ independently
of the coupling---on the same chip):

\begin{center}\small
\resizebox{\textwidth}{!}{%
\begin{tabular}{lccccc}
\toprule
device & $\kappa$ [cm$^{-1}$] & $\gamma$ [cm$^{-1}$] & regime &
predicted $t_c$ [cm] & predicted floor [cm$^{-1}$] (D3: window period)\\
\midrule
D1 & 0.10 & 0.60 & broken ($\gamma/2\kappa=3$) & 5.10 & $-0.583$\\
D2 & 0.1125 & 0.90 & broken, deeper ($\gamma/2\kappa=4$) & 3.98 & $-0.886$\\
D3 & 0.50 & 0.30 & unbroken & 1.55 (first window) &
period $6.59\,$cm\\
D4 & 0.30 & 0.60 & \emph{at} the EP ($\gamma/2\kappa=1$) &
2.36 ($=1/\kappa\sqrt2$) & $-0.300$ ($=-a$ asymptote, $\tilde\kappa=0$)\\
\bottomrule
\end{tabular}%
}
\end{center}

\noindent For each family, a ladder of chip lengths
$z=0.5,1.0,\dots,10\,$cm (20 samples; Klauck-style
varying-length sets). At each length, launch classical laser light
into each guide in turn and record both output intensities relative
to the straight reference guide: four numbers per length assemble
$\hat\Gamma(z)$. The required loss rates ($2.6$--$7.8\,$dB/cm power
loss) are set by the tunable-bending-loss modulation
technique~\cite{Eichelkraut2014}, which reports radiation losses of
almost arbitrary amount; we note that the highest loss demonstrated on
this same platform in~\cite{Klauck2019} is ${\approx}3\,$dB/cm, so D2's
requirement sits a factor ${\sim}2.6$ above the published
same-platform record. D2, the deepest working device,
keeps its whole measurement---onset and floor---within $6\,$cm chips. A single-photon version
(heralded pairs, as in~\cite{Klauck2019}) would reproduce the numbers
below at $10^4$ detected photons per launch; entangled-pair
extensions~\cite{Tang2016,Ehrhardt2022} address the
conditioning-signal observable treated elsewhere in this program and
are out of scope here.

\paragraph{Estimators and error budget (simulated).} The
measurement-chain simulation models per-sample fabrication disorder
($2\%$ on $\kappa$, $a$), per-launch transmission-calibration error
($3\%$), photon shot noise, and a $10^{-4}$ dark/crosstalk floor,
with $400$ Monte Carlo realizations per configuration:
\begin{itemize}\setlength\itemsep{0pt}
\item \emph{$t_c$ is calibration-invariant by construction:} the onset
is estimated from the zero crossing of $\det$ of the
\emph{column-normalized} transfer matrix, which is exactly invariant
under per-launch intensity rescalings---an algebraic identity, and the
simulated calibration-only spread is correspondingly zero.
Precision at $10^4$ photons/launch: $\pm0.19\,$cm (D1),
$\pm0.28\,$cm (D2); at $10^6$ with all systematics, $\pm0.09$ and
$\pm0.06\,$cm (Klauck's HOM-dip benchmark:
$0.05\,$cm~\cite{Klauck2019}). The dark floor is the one systematic
this estimator does \emph{not} absorb: it biases the crossing by
$+0.05\,$cm (D1) and $+0.12\,$cm (D2), and must be subtracted. The
bias grows with $\gamma/2\kappa$ because the surviving diagonal at
$t_c$ weakens; by $\gamma/2\kappa\approx5$--$6$ the determinant no
longer cleanly crosses zero at this dark level and the estimator
degrades beyond usefulness, its bias and spread exceeding the signal,
which is what sets D2's depth.
\item \emph{The floor requires the global fit:} pointwise
finite-difference generators are fragile where the survival
probability approaches the dark floor (documented in the script's
null control). Fitting the closed-form model to the full ladder, with
per-column calibration nuisances, recovers the floor to $\pm3.4\%$
(D1) and $\pm4.1\%$ (D2)
for both broken devices---$29\sigma$ (D1) and $24\sigma$ (D2) below
zero---already at $10^4$ photons/launch. This is the significance of
the floor \emph{parameter} under the closed-form model, not a
model-independent test of permanence.
\item \emph{Dissociation:} on identical data and analysis, the
late-window ($z\ge5\,$cm) conditional backflow is $0.018$ (D2, whose
onset $t_c=3.98\,$cm precedes the window---the clean broken-phase
instance, dropping to $0.004$ at $10^6$ photons) and $0.015$ (D1, whose
$t_c=5.10\,$cm sits at the window edge, so this retains the
theory-predicted onset-neighborhood tail---$\approx0.01$ noiseless,
$0.009$ as simulated at $10^{6}$, the dark floor slightly suppressing it---%
rather than pure noise) versus $0.49$ for D3
(unbroken), while the fitted floor is $\ge24\sigma$ negative for D1/D2:
the dissociation is resolved with more than an order of magnitude of
margin.
\item \emph{Unbroken control:} D3's first window position is
recovered to $\pm0.03\,$cm and the window period to $\pm2.3\%$,
confirming the recurrent (non-terminal) signature. (The
exceptional-point device D4 is not in the simulated budget; continuity
of $t_c$ through the EP is an analytic property of the closed form, not
a simulated result here.)
\item \emph{Two dimensionless points:} $\kappa\,t_c$ and
$\mathrm{floor}/\kappa$ depend on $\gamma/2\kappa$ alone, so D1 and D2
test the closed form at $\gamma/2\kappa=3$ and $4$
($\kappa\,t_c=0.510$ and $0.447$) rather than re-testing one
prediction at two overall scales.
\end{itemize}

\paragraph{What a positive result establishes.} These would be the
first process-level (transition-matrix) observables of $\PT$ dynamics:
they test the identification of broken-$\PT$ evolution with
survival-conditioned dynamics of a dilated stochastic process---in
particular that the broken phase is \emph{backflow-Markovian at
late lengths yet permanently indivisible}, a prediction with no analogue in the
quantum-state-based analyses of~\cite{Xiao2019,KAU2017}. A null
result---most directly, a genuine \emph{restoration} of divisibility
($\det\hat\Gamma$ crossing back positive) in a resolvable window just
past $t_c$, a model-independent signal; or a systematic late-window
lack of fit of the closed form
(e.g.\ $L_{12}$ relaxing toward zero once its asymptote is freed from
$(a,\kappa)$ in an extended fit)---would falsify the closed forms and,
with the device Hamiltonian independently constrained by the on-chip
reference structures, the dilation reading for this system. The
restoration test must be read only where the divisor determinant is
resolvable: since $\det\hat\Gamma\to0^-$ as $z\to\infty$ (survival, and
with it every intensity, vanishing), a naive late-$z$ sign reading is
powerless---at the $10^{-4}$ dark level and $10^{4}$ photons per launch
the measured late-ladder
determinant flips positive in $79\%$ of realizations even when the closed
forms hold exactly (falling to $9\%$ at $10^{6}$)---so a genuine restoration must clear the measured
det noise-plus-dark envelope in the window just past $t_c$, not merely
change sign at large $z$. The lack-of-fit falsifier is informative
already at $10^4$ detected photons per launch (an afternoon of data on
this platform); the model-independent restoration test additionally
requires photons enough to resolve $\det\hat\Gamma$ above the dark
floor across that window.

\medskip
\noindent\textbf{Provenance.} Prepared with the large language model
Claude Fable~5 (Anthropic), directed by the author; simulation and
error budget reproducible from
\texttt{scripts/t2\_measurement\_chain.py} in the public repository
\url{https://github.com/embeddetech/p02-measurable-ghosts}
(also archived in the program hub,
\url{https://github.com/embeddetech/ghosts-program}).

\begin{thebibliography}{9}\small
\bibitem{indivisible}
J.~Torkelson,
\emph{Indivisible stochastic processes for non-Hermitian and
higher-derivative quantum systems: detailed balance, broken $\PT$
symmetry, and the explosion of ghosts},
paper~1 of this program, repository
\url{https://github.com/embeddetech/p01-indivisible-ghosts}
(2026).
\bibitem{Klauck2019} F.~Klauck \emph{et al.},
\emph{Observation of $\PT$-symmetric quantum interference},
Nat.\ Photonics \textbf{13}, 883 (2019),
\href{https://arxiv.org/abs/1904.08135}{arXiv:1904.08135}.
\bibitem{Eichelkraut2014} T.~Eichelkraut \emph{et al.},
\emph{Radiation-loss management in modulated waveguides},
Opt.\ Lett.\ \textbf{39}, 6831 (2014).
\bibitem{Xiao2019} L.~Xiao \emph{et al.},
\emph{Observation of critical phenomena in parity-time-symmetric
quantum dynamics}, Phys.\ Rev.\ Lett.\ \textbf{123}, 230401 (2019).
\bibitem{Bernardes2015} N.~K.~Bernardes \emph{et al.},
\emph{Experimental observation of weak non-Markovianity},
Sci.\ Rep.\ \textbf{5}, 17520 (2015).
\bibitem{Tang2016} J.-S.~Tang \emph{et al.},
\emph{Experimental investigation of the no-signalling principle in
parity--time symmetric theory using an open quantum system},
Nat.\ Photonics \textbf{10}, 642 (2016).
\bibitem{Ehrhardt2022} M.~Ehrhardt, M.~Heinrich, and A.~Szameit,
\emph{Observation-dependent suppression and enhancement of two-photon
coincidences by tailored losses},
Nat.\ Photonics \textbf{16}, 191 (2022).
\bibitem{KAU2017} K.~Kawabata, Y.~Ashida, and M.~Ueda,
\emph{Information retrieval and criticality in parity-time-symmetric
systems}, Phys.\ Rev.\ Lett.\ \textbf{119}, 190401 (2017).
\end{thebibliography}

\end{document}
